\section{Maple Tree 后端：\texttt{vma\_maple.c}}
\label{sec:vma-maple}

Linux 6.1（2022 年 12 月）用 Maple Tree 取代了沿用 20 年的红黑树来管理 VMA。
我们实现了一个教学级的 B+ tree 变体，保留 Maple Tree 的核心设计理念。

\begin{whybox}
为什么 Linux 从红黑树换到 Maple Tree？

\begin{enumerate}
  \item \textbf{缓存友好}：B-tree 节点包含多个键值对（扇出度高），
        一次缓存行加载可比较多个键，减少缓存 miss。
        红黑树每个节点只有一个键，每次比较都可能 cache miss。
  \item \textbf{RCU 安全}：Maple Tree 天然支持 Read-Copy-Update，
        读者无需加锁即可遍历。红黑树的旋转在 RCU 下很难正确实现。
  \item \textbf{区间原生支持}：以 \texttt{[start, end)} 区间为键，
        不像红黑树只按 \texttt{start} 排序后再检查 \texttt{end}。
  \item \textbf{减少指针追逐}：紧凑数组存储，指针数更少。
\end{enumerate}
\end{whybox}

\subsection{节点结构}

每个节点包含最多 9 个分割键（pivot）和 10 个子节点/数据槽位：

\codefile{src/kernel/mm/vma\_maple.c}
\begin{lstlisting}[style=cstyle]
#define MAPLE_NODE_SLOTS    10
#define MAPLE_MAX_PIVOTS    (MAPLE_NODE_SLOTS - 1)  /* 9 */

typedef struct maple_node {
    uint32_t      pivots[MAPLE_MAX_PIVOTS];
    uint8_t       nr_entries;
    uint8_t       is_leaf;
    maple_node_t* parent;
    union {
        maple_node_t* children[MAPLE_NODE_SLOTS]; /* 内部节点 */
        vm_area_t*    slots[MAPLE_NODE_SLOTS];    /* 叶节点 */
    };
} maple_node_t;
\end{lstlisting}

\begin{figure}[H]
  \centering
  % figures/ch11/fig08-maple-node.tex — auto-extracted from chapter source
\begin{tikzpicture}[
  cell/.style={draw, minimum height=0.7cm, minimum width=1.0cm,
               font=\ttfamily\scriptsize, anchor=west},
  arr/.style={-{Stealth[length=4pt]}, thick},
]
  % Internal node
  \node[font=\small\bfseries] at (-1.5, 0.3) {内部节点:};
  \node[cell, fill=orange!15] (p0) at (0, 0) {P0};
  \node[cell, fill=orange!15] (p1) at (1.0, 0) {P1};
  \node[cell, fill=orange!15] (p2) at (2.0, 0) {$\cdots$};
  \node[cell, fill=orange!15] (p8) at (3.0, 0) {P8};

  \node[cell, fill=blue!10] (c0) at (0, -1.0) {C0};
  \node[cell, fill=blue!10] (c1) at (1.0, -1.0) {C1};
  \node[cell, fill=blue!10] (c2) at (2.0, -1.0) {$\cdots$};
  \node[cell, fill=blue!10] (c9) at (3.0, -1.0) {C9};

  \node[font=\scriptsize\color{gray}] at (4.5, 0) {9 pivots};
  \node[font=\scriptsize\color{gray}] at (4.5, -1.0) {10 children};

  % Leaf node
  \node[font=\small\bfseries] at (-1.5, -2.5) {叶节点:};
  \node[cell, fill=orange!15] at (0, -2.8) {P0};
  \node[cell, fill=orange!15] at (1.0, -2.8) {P1};
  \node[cell, fill=orange!15] at (2.0, -2.8) {$\cdots$};
  \node[cell, fill=orange!15] at (3.0, -2.8) {P8};

  \node[cell, fill=green!12] at (0, -3.8) {VMA};
  \node[cell, fill=green!12] at (1.0, -3.8) {VMA};
  \node[cell, fill=green!12] at (2.0, -3.8) {$\cdots$};
  \node[cell, fill=green!12] at (3.0, -3.8) {VMA};

  \node[font=\scriptsize\color{gray}] at (4.5, -2.8) {9 pivots};
  \node[font=\scriptsize\color{gray}] at (4.5, -3.8) {10 VMA ptrs};
\end{tikzpicture}

  \caption{Maple Tree 节点布局——内部节点存子指针，叶节点存 VMA 指针}
  \label{fig:maple-node}
\end{figure}

\begin{whybox}
为什么扇出度选 10？

Linux 的 \texttt{maple\_range\_64} 有 16 个 slot，每节点 256 字节（4 个缓存行）。
我们的 32 位系统中，10 个 pivot + 10 个 child + 元数据 $\approx$ 96 字节（$<$ 2 个缓存行）。
扇出度 10 意味着 256 个 VMA 只需树高 3（$10^3 = 1000 > 256$），
每次查找最多 3 次节点访问，每次访问时同一节点的 pivot 在同一缓存行中。
\end{whybox}

\subsection{查找：从根下降到叶}

与红黑树不同，Maple Tree 的查找直接沿着 B+ tree 结构下降到叶节点，
然后在叶中线性扫描 $\leq 10$ 个条目：

\codefile{src/kernel/mm/vma\_maple.c}
\begin{lstlisting}[style=cstyle]
static const vm_area_t* maple_find(uint32_t addr)
{
    maple_node_t* node = g_maple_root;
    while (!node->is_leaf) {
        unsigned idx = internal_find_child(node, addr);
        node = node->children[idx];
    }
    for (unsigned i = 0; i < node->nr_entries; i++) {
        vm_area_t* vma = node->slots[i];
        if (vma && addr >= vma->start && addr < vma->end)
            return vma;
    }
    return NULL;
}
\end{lstlisting}

\subsection{插入：分裂上提}

当叶节点满了（10 个条目），需要分裂：

\begin{enumerate}
  \item 分配新叶节点，把当前叶的后半部分移到新叶
  \item 把中间键上提到父节点
  \item 如果父节点也满了，递归分裂——这是 B-tree 增高的唯一方式
\end{enumerate}

\begin{figure}[H]
  \centering
  % figures/ch11/fig09-maple-split.tex — auto-extracted from chapter source
\begin{tikzpicture}[
  cell/.style={draw, minimum height=0.6cm, minimum width=0.6cm,
               font=\ttfamily\scriptsize, anchor=west},
  arr/.style={-{Stealth[length=4pt]}, very thick, red!60!black},
]
  % Before split
  \node[font=\small\bfseries] at (-1.5, 0) {分裂前:};
  \foreach \i/\c in {0/K0,1/K1,2/K2,3/K3,4/K4,5/K5,6/K6,7/K7,8/K8,9/K9} {
    \node[cell, fill=green!12] at (\i*0.7, 0) {\c};
  }
  \node[font=\scriptsize\color{red}] at (8.2, 0) {满！};

  % After split
  \node[font=\small\bfseries] at (-1.5, -1.5) {分裂后:};
  \foreach \i/\c in {0/K0,1/K1,2/K2,3/K3,4/K4} {
    \node[cell, fill=green!12] at (\i*0.7, -1.5) {\c};
  }
  \foreach \i/\c in {0/K5,1/K6,2/K7,3/K8,4/K9} {
    \node[cell, fill=blue!12] at (5.0+\i*0.7, -1.5) {\c};
  }

  \node[cell, fill=orange!20] (up) at (3.5, -2.8) {K5};
  \node[font=\scriptsize] at (4.5, -2.8) {$\uparrow$ 上提到父节点};

  \draw[arr] (3.5, -1.1) -- (3.5, -2.5);
\end{tikzpicture}

  \caption{Maple Tree 叶节点分裂——后半部分移到新节点，中间键上提}
  \label{fig:maple-split}
\end{figure}

\subsection{与红黑树的对比}

\begin{designbox}
两种后端的适用场景：
\begin{itemize}
  \item \textbf{红黑树}：通用自平衡 BST，实现经典，适合教学。
        每个节点只存一个 VMA，指针追逐多，但代码逻辑清晰。
  \item \textbf{Maple Tree}：面向现代 CPU 缓存层次优化。
        每次节点访问比较多个键（同一缓存行），树更矮，总缓存 miss 更少。
        代价是分裂/合并逻辑比旋转更复杂。
\end{itemize}
在 VMA 数量 $< 256$ 的内核启动阶段，三种后端的性能差异可忽略。
选择哪个主要看代码风格偏好和教学目标。
\end{designbox}

% ============================================================================

